\section*{第 \ref{ch:b2:biomolecules} 章——生物分子：氨基酸、糖、脂质与核苷酸}

\begin{solution}{exo:b2:biomolecules:1}
\begin{itemize}
  \item 在 pH 3 至 9 之间，两性离子 \ce{H3N+-CH(CH3)-COO-}。
  \item 在 pH 1 时，阳离子 \ce{H3N+-CH(CH3)-COOH}。
  \item 在 pH 12 时，阴离子 \ce{H2N-CH(CH3)-COO-}。
\end{itemize}
\end{solution}

\begin{solution}{exo:b2:biomolecules:2}
甘氨酸：$(2.35 + 9.78)/2 \approx 6.1$。丙氨酸：$(2.35 + 9.87)/2 \approx 6.1$。
% ledger: pka:glycine1, pka:glycine2, pka:alanine1, pka:alanine2
\end{solution}

\begin{solution}{exo:b2:biomolecules:3}
\ce{H2N-CH(CH3)-CO-NH-CH2-COOH}：丙氨酸提供 N 端，甘氨酸提供 C 端。Gly--Ala，
\ce{H2N-CH2-CO-NH-CH(CH3)-COOH}为一种构造异构体：游离氨基位于甘氨酸上。
\end{solution}

\begin{solution}{exo:b2:biomolecules:4}
\ce{NH2} 位于左侧，为 L 构型。优先顺序为 \ce{NH2} > \ce{COOH} > \ce{CH2OH} > \ce{H}：L-丝氨酸为
（\textit{S}）。
\end{solution}

\begin{solution}{exo:b2:biomolecules:5}
基团：N 端铵基（约 9）、赖氨酸侧链（10.7）、天冬氨酸侧链（3.9），
C 端羧基（约 2）。pH 1 时：$+1 +1 + 0 + 0 = +2$。pH 7 时：$+1 +1 -1 -1 = 0$。pH 12 时：
$0 + 0 - 1 - 1 = -2$（赖氨酸侧链比 pH 12 低 1.3 个单位，主要呈中性）。
\end{solution}

\begin{solution}{exo:b2:biomolecules:6}
氯化亚砜加热还会攻击其他基团，尤其是活性很强的酰氯，
而氨基酸的酰氯很容易失去构型（$\alpha$ 氢变得呈酸性），使肽
部分发生外消旋化。偶联剂可在室温、温和条件下活化羧酸。
\end{solution}

\begin{solution}{exo:b2:biomolecules:7}
O5 与 C2 成键：由 C2、C3、C4、C5 和 O 组成的环，共五个原子（呋喃糖）。在哈沃斯投影式中，环上的 O
位于后方，C2 在右侧，连有 \ce{CH2OH}（C1）朝下，\ce{OH} 朝上（$\beta$）；C3 的 \ce{OH} 朝上（
在费歇尔投影式中居左），C4 的 \ce{OH} 朝下，C5 连有 \ce{CH2OH}（C6）朝上。
\end{solution}

\begin{solution}{exo:b2:biomolecules:8}
$x_\alpha = (80.2 - 52.8)/(150.7 - 52.8) = 27.4/97.9 \approx 0.28$：28~\% $\alpha$，72~\% $\beta$。
\end{solution}

\begin{solution}{exo:b2:biomolecules:9}
麦芽糖在第二个葡萄糖单元上保留游离的半缩醛，可开环生成醛基，具有还原性。在蔗糖中
两个异头碳都参与糖苷键的形成，不存在半缩醛，因而不具还原性。稀酸
可水解糖苷键（一种缩醛），释放出葡萄糖和果糖，两者均具还原性。
\end{solution}

\begin{solution}{exo:b2:biomolecules:10}
三油酸甘油酯，\qty{885.45}{g/mol}，消耗三个 \ce{KOH}（\qty{56.105}{g/mol}）：$3 \times 56.105/885.45 =
0.190$~g/g，即皂化值为 190。链更短时，
同样三个酯基对应的摩尔质量更小：每克脂肪的摩尔数更多，所需 KOH 也更多。
% ledger: aw:C, aw:H, aw:O, aw:K
\end{solution}

\begin{solution}{exo:b2:biomolecules:11}
将甘氨酸的氨基以 \textit{tert}-丁氧羰基氨基甲酸酯形式保护，丙氨酸的羧基以
甲酯形式保护。以碳二亚胺偶联。用酸脱去氨基甲酸酯，用稀碱脱去
酯（条件正交，两端均可单独释放）：Gly--Ala。
\end{solution}

\begin{solution}{exo:b2:biomolecules:12}
5$'$-ACGCAT-3$'$（反向平行读取：A 与 T 配对，G 与 C 配对）。配对：A--T、T--A、G--C、C--G、G--C、T--A：
$3 \times 2 + 3 \times 3 = 15$ 个氢键。
\end{solution}

\begin{solution}{pb:b2:biomolecules:1}
\textbf{1.} L-天冬氨酸、L-苯丙氨酸和甲醇。
\textbf{2.} 两者中，\ce{COOH} 位于上方，侧链位于下方（\ce{CH2COOH}、\ce{CH2C6H5}），
\ce{NH2} 位于左侧。
\textbf{3.} 两者均为（\textit{S}）。
\textbf{4.} 两个手性碳原子：共四种立体异构体。
\textbf{5.} \ce{H2N-CH(CH2COOH)-CO-NH-CH(CH2C6H5)-COOCH3}：肽键连接两个
残基；酯基为 C 端的 \ce{COOCH3}。
\textbf{6.} 味觉受体是手性的：它们只与一种立体异构体结合，而不与其对映体
或非对映异构体结合，正如引言所述。
\textbf{7.} N 端的 \ce{NH3+} 与天冬氨酸侧链上的 \ce{COOH}；C 端的羧基
为甲酯，不含酸性氢。
\textbf{8.} pH 2 时：\ce{NH3+}（$+1$）、侧链 \ce{COOH}（0）：$+1$。
\textbf{9.} pH 7 时：$+1 - 1 = 0$。
\textbf{10.} pH 12 时：$0 - 1 = -1$。
\textbf{11.} 在发生变化的两个基团的 $\mathrm pK_a$ 之间：$(3.9 + 9.9)/2 \approx 6.9$。
\textbf{12.} 在二肽中，氨基所连的碳上带有酰胺基而非羧酸根，
相邻基团和电荷随之改变，$\mathrm pK_a$ 也随之变化（肽的 N 端铵基
明显比游离氨基酸的铵基更酸）。该模型只给出形状，而非
精确数值。
\textbf{13.} 羧酸与胺先形成盐；加热会得到混合物（Asp--Asp、更长的链、
错误的羧基），并使手性碳原子发生外消旋化。
\textbf{14.} 其氨基和侧链羧基。
\textbf{15.} 它将羧酸的 \ce{OH} 转变为良好的离去基团，形成活化衍生物，
胺进攻该衍生物即可在室温下发生反应（酰基取代）。
\textbf{16.} 该 $\beta$-二肽通过侧链连接：是阿斯巴甜的异构体，而非阿斯巴甜本身。
\textbf{17.} 氨基以苄氧羰基氨基甲酸酯形式保护，侧链以苄酯形式保护，两者
可在钯上经氢解一并脱去，甲酯保持不变（与其
正交）。
\textbf{18.} 活化的羧酸具有酸性的 $\alpha$ 氢，碱会将其除去，使手性碳原子
发生外消旋化。
\textbf{19.} 保护天冬氨酸（N 上为氨基甲酸酯，侧链为苄酯）；活化游离的
$\alpha$-\ce{COOH} 用偶联剂活化；加入 L-苯丙氨酸甲酯；进行氢解。
\textbf{20.} 甲酯与肽键；酯基更易反应，因为酰胺的反应活性低得多。
\textbf{21.} 二肽 Asp--Phe（游离酸）和甲醇。
\textbf{22.} 甜味分子被消耗了：其水解产物不能代替它。
\textbf{23.} $14 \times 12.011 + 18 \times 1.008 + 2 \times 14.007 + 5 \times 15.999 =
\qty{294.307}{g/mol}$。
\textbf{24.} $0.180/294.307 = \qty{6.116e-4}{mol}$，\qty{0.612}{mmol}。
\textbf{25.} 每个阿斯巴甜对应一个甲醇，因此甲醇的质量为 $0.6116 \times 32.042 = \boldsymbol{\approx \qty{19.6}{mg}}$
。
\end{solution}
